\section{线性表示与整体李群}

我们所谈论的这些工作，没有一项坦率地触及整体李群的定义与研究问题。沿这条道路的最初几步要追溯到 H. 外尔。他从两个直至当时仍独立发展的理论中得到启发：E. 嘉当的复半单李代数线性表示理论，以及弗罗贝尼乌斯的有限群线性表示理论——后者刚由 I. 舒尔利用胡尔维茨的一个想法移植到正交群上。胡尔维茨曾证明{[}168{]}，通过把有限群上的平均运算换成关于一个不变测度的积分，可以为正交群或酉群构造不变量。他还注意到，把这个方法应用于酉群时，就得到一般线性群的不变量，这是“酉技巧”的第一个例子。1924 年，I. 舒尔{[}279 e{]}用这一程序，通过构造一个正定的非退化不变埃尔米特型，证明了正交群 O(n) 与酉群 U(n) 的表示的完全可约性；他由此借助“酉技巧”推出 O(n, C) 与 SL(n, C) 的全纯表示的完全可约性，为 O(n) 与 U(n) 的特征标建立了正交关系，并确定了 O(n) 的特征标。H. 外尔还把这个方法推广到复半单李代数{[}331 b{]}。给定这样一个代数 g，他证明它具有一个“实紧形式”（这等于说它来自 R 上的一个代数 \(g_{0}\)——其伴随群 \(G_{0}\) 是紧的——把纯量从 R 扩张到 C 而得）。此外他还证明 \(G_{0}\) 的基本群是有限的，因而泛覆盖\(^{20}\)（\(G_{0}\) 的）是紧的。他由此，通过对舒尔程序的一个适当改造，推出 g 的表示的完全可约性，并且还用整体手段给出了 g 的表示的特征标的确定。在一封给 I. 舒尔的信{[}331 a{]}中，H. 外尔概述了舒尔当时尚不知道的嘉当的结果（参见{[}279 e{]}，p.~299, 页脚注），并比较了两种观点：嘉当的方法给出李代数 g 的单连通群的所有全纯表示；在正交群的情形，还得到了一个二叶覆盖（后来称为旋量群）的表示，这是舒尔未曾注意到的；另一方面，舒尔的方法的优点在于能证明完全可约性并显式地给出特征标。

在 H. 外尔的工作之后，E. 嘉当在对对称空间与李群的研究中采取了公开的整体观点。正是这一观点成为他 1930 年对“有限连续”群理论的阐述（{[}52 a{]}，v. I₂, pp.~1165-1225）的基础。其中特别可以找到第三基本定理的整体变体的第一个证明（具有给定李代数的李群的存在性）；嘉当还证明，实李群的每个闭子群都是李群，这推广了 J. 冯·诺伊曼关于线性群的闭子群的一个结果{[}324 b{]}。在这篇论文中，冯·诺伊曼还证明了复半单群的每个连续表示都是实解析的。

在这项工作之后，通常意义（即在 R 或 C 上有限维）下的李群理论其主要轮廓已大体固定。对它的第一个详细阐述由庞特里亚金在他的拓扑群著作{[}253{]}中给出；他在书中坚持一个仍相当接近李的观点，但仔细地区分了局部与整体。随后是谢瓦莱的书{[}62 d{]}，其中还包含对解析流形理论与外微分演算的第一次系统讨论；李的“无穷小变换”在那里表现为向量场，李群 G 的李代数被等同于在 G 下左不变的向量场所成的空间。他把“群芽”方面与“变换群”方面放在了一边。

